% Cote 119: offprints, reports, courses, CV and letters (typed and printed part): archive cote 119, site page 175 (PDF page 176).
% typescript. Draft compared: none.
% AI-transcribed by Claude Opus 5.5 (claude-opus-5-5, adaptive thinking, effort high), Message Batches API, batch r119_full (msgbatch_01DwZNm5Y7YsKBWyAN4sFs8w), 2026-11-13; not yet checked by eye.
\scanpage{119/175}
\mspage{- 14 -}

\uline{Démonstration .}

Soit \( r < \rho \). On va trouver une série entière qui converge normale-%
ment vers \( f(Z) \) pour \( |Z| \leq r \) . Cette série sera indépendante de \( r \) en vertu
de l'unicité du développement en série entière d'une fonction au voisinage de
\( Z_0 \). Le théorème en résultera donc.

Choisissons \( r_0 \) tel que \( r < r_0 < \rho \) . Appliquons le théorème 5.4.
en prenant pour \( \gamma \) le cercle de centre \( Z_0 \) et de rayon \( r_0 \). Pour \( |Z-Z_0| \leq r \)
on a : \( f(Z) = \frac{1}{2 \pi i} \int_{\gamma} \frac{f(t)\, dt}{t - Z} \)
Comme \( |Z-Z_0| < |t-Z_0| \) on peut développer \( \frac{1}{t - Z} \) :
\[ \frac{1}{t - Z} = \frac{1}{t - Z_0} \;.\; \frac{1}{1 - \frac{Z-Z_0}{t-Z_0}} = \frac{1}{t - Z_0} \Bigl( 1 + \frac{Z - Z_0}{t - Z_0} + \Bigl( \frac{Z - Z_0}{t - Z_0} \Bigr)^2 + \dots \Bigr) \]

On a donc :
\[ f(Z) = \frac{1}{2 \pi \, i} \int_{\gamma} \sum_{n \geq 0} \frac{(Z-Z_0)^n}{(t-Z_0)^{n+1}} f(t)\, dt \]

La série converge normalement pour \( |Z - Z_0| \leq r \) et \( |t - Z_0| = r_0 \) on
peut donc l'intégrer terme à terme et on obtient une série normalement conver-%
gente pour \( |Z - Z_0| \leq r \) . (Rappeler la justification de ce fait constitue
un bon exercice).

On obtient donc :
\[ (5.5.1.) \qquad f(Z) = \sum_{n \geq 0} a_n \, (Z - Z_0)^n \]
où les coefficients \( a_n \) sont donnés par les intégrales
\[ (5.5.2.) \qquad a_n = \frac{1}{2 \pi i} \int_{|t-Z_0| = r_0} \frac{f(t)\, dt}{(t-Z_0)^{n+1}} \]

Le théorème est donc démontré.

(5.6.) \uline{Remarque.}

On a obtenu plus que le résultat attendu puisque non seulement on
a prouvé que toute fonction holomorphe est analytique mais on a montré de plus

\margin{175}
